% Source: https://artofproblemsolving.com/community/c3h1987760_counting_problem
% Stable identity: aops-topic-1987760-q01
An arrangement is a sequence of alternating blocks (runs) of $A$'s and $B$'s, and the condition says every run has length $1$ or $2$. Write $4$ as an ordered sum of $1$'s and $2$'s: with $k$ parts there are $a_k$ ways, where $a_2=1$ $(2+2)$, $a_3=3$ $(1+1+2$ in some order$)$ and $a_4=1$ $(1+1+1+1)$.

If the $A$'s form $k$ runs and the $B$'s form $j$ runs, the runs alternate, so $|k-j|\le1$. When $k=j$ the arrangement can start with either letter ($2$ ways); when $|k-j|=1$ it must start with the letter that has more runs ($1$ way). Hence the count is
\[
2\left(a_2^2+a_3^2+a_4^2\right)+2\left(a_2a_3+a_3a_4\right)=2(1+9+1)+2(3+3)=22+12=\boxed{34}.
\]
